% Burgers equation with viscous/artificial-viscosity shock
% Compiled from the discussion

\documentclass{mooiman_memo}


\usepackage{listings}

\memoSubject{Burgers Equation with Viscous and Artificial-Viscosity Shock}
\memoFrom{Jan Mooiman}
\memoDate{\today~\currenttime}
\memoTo{}


\begin{document}

\maketitle

\section{Problem}

Consider the inviscid Burgers equation
%
\begin{equation}
    \frac{\partial u}{\partial t}
    +u\frac{\partial u}{\partial x}=0,
    \qquad 0<x<L_x,
\end{equation}
%
with boundary conditions
%
\begin{equation}
    u(0,t)=a,
    \qquad
    u(L_x,t)=b.
\end{equation}
%
For the time-dependent problem, an initial condition
\begin{equation}
    u(x,0)=u_0(x)
\end{equation}
is also required.
%
The equation can be written in conservation form as
%
\begin{equation}
    u_t+\left(\frac{u^2}{2}\right)_x=0.
\end{equation}
%
\section{Characteristics}
%
Along a characteristic,
%
\begin{equation}
    \frac{dx}{dt}=u,
    \qquad
    \frac{du}{dt}=0.
\end{equation}
%
Thus a constant value $u_0$ propagates with speed $u_0$:
%
\begin{equation}
    x=x_0+u_0t.
\end{equation}
5
\section{Shock solution}
%
For a shock connecting a left state $a$ to a right state $b$, the
Rankine--Hugoniot shock speed is
%
\begin{equation}
    s=
    \frac{f(a)-f(b)}{a-b},
    \qquad
    f(u)=\frac{u^2}{2}.
\end{equation}
%
Consequently,
%
\begin{equation}
    \boxed{
    s=\frac{a+b}{2}.
    }
\end{equation}
%
The shock position is therefore
%
\begin{equation}
    x_s(t)=x_s(0)+\frac{a+b}{2}t.
\end{equation}
%
A stationary shock requires
%
\begin{equation}
    \boxed{a+b=0}
\end{equation}
%
or
%
\begin{equation}
    \boxed{b=-a}.
\end{equation}
%
For this case, the inviscid stationary solution is
%
\begin{equation}
    u(x)=
    \begin{cases}
        a, & 0\leq x<x_s,\\
        b, & x_s<x\leq L_x.
    \end{cases}
\end{equation}
%
The inviscid equation alone does not determine $x_s$ when $b=-a$;
any shock location inside the domain satisfies the stationary
Rankine--Hugoniot condition.

For the frequently used case
%
\begin{equation}
    a=1,\qquad b=-1,
\end{equation}
5
the shock speed is
%
\begin{equation}
    s=0.
\end{equation}
%
\section{Viscous Burgers equation}
5
Introduce a constant viscosity,
%
\begin{equation}
    \boxed{
    u_t+u u_x=\nu_{\rm art}u_{xx}.
    }
\end{equation}
%
For a stationary shock,
%
\begin{equation}
    u u_x=\nu_{\rm art}u_{xx}.
\end{equation}
%
Integrating once gives
%
\begin{equation}
    \nu_{\rm art}u_x
    =
    \frac{u^2}{2}+C.
\end{equation}
%
For the symmetric case $a>0$, $b=-a$,
%
\begin{equation}
    C=-\frac{a^2}{2},
\end{equation}
%
so
%
\begin{equation}
    \nu_{\rm art}u_x
    =
    \frac12(a^2-u^2).
\end{equation}
%
The viscous shock solution connecting $a$ on the left to $-a$ on the
right is
%
\begin{equation}
    \boxed{
    u(x)=
    -a\tanh\left[
        \frac{a}{2\nu_{\rm art}}(x-x_s)
    \right].
    }
\end{equation}
%
For $a=1$,
%
\begin{equation}
    \boxed{
    u(x)=
    -\tanh\left[
        \frac{x-x_s}{2\nu_{\rm art}}
    \right].
    }
\end{equation}
%
The characteristic shock length scale is approximately
%
\begin{equation}
    \ell_s\sim 2\nu_{\rm art}
\end{equation}
%
for $a=1$, and more generally
%
\begin{equation}
    \ell_s\sim\frac{2\nu_{\rm art}}{a}.
\end{equation}
%
\section{General viscous travelling shock}
%
For arbitrary left and right states $a>b$, the travelling viscous
profile is approximately
%
\begin{equation}
    \boxed{
    u(x,t)=
    \frac{a+b}{2}
    -
    \frac{a-b}{2}
    \tanh\left[
        \frac{a-b}{4\nu_{\rm art}}
        \left(x-x_s(t)\right)
    \right],
    }
\end{equation}
%
where
%
\begin{equation}
    x_s(t)=x_s(0)+\frac{a+b}{2}t.
\end{equation}
%
For $a=1$, $b=-1$, this reduces to
%
\begin{equation}
    u(x)=
    -\tanh\left[
        \frac{x-x_s}{2\nu_{\rm art}}
    \right].
\end{equation}
%
\section{Artificial-viscosity discretisation}
%
The artificial-viscosity matrix contribution used in the numerical
scheme is
%
\begin{equation}
    c_\psi
    \begin{bmatrix}
        -1 & 2 & -1
    \end{bmatrix}.
\end{equation}
%
Its action on the solution is
%
\begin{equation}
    -c_\psi u_{i-1}
    +2c_\psi u_i
    -c_\psi u_{i+1}.
\end{equation}
%
Using the Taylor expansion
%
\begin{align}
    u_{i-1}-2u_i+u_{i+1}
    &=
    \Delta x^2u_{xx}
    +\frac{\Delta x^4}{12}u_{xxxx}
    +O(\Delta x^6),
\end{align}
%
the artificial-viscosity contribution becomes
%
\begin{equation}
    \boxed{
    -c_\psi\Delta x^2u_{xx}
    -\frac{c_\psi\Delta x^4}{12}u_{xxxx}
    +O(\Delta x^6).
    }
\end{equation}
%
If the matrix is normalized so that this term represents a standard
diffusion term, the leading coefficient is associated with
%
\begin{equation}
    \nu_{\rm art}\sim c_\psi\Delta x^2.
\end{equation}
%
However, the complete scheme contains an additional right-hand-side
term and therefore this identification should not be made without
considering the complete discrete equation.
%
\section{Mass matrix used in the scheme}
%
The mass stencil is
%
\begin{equation}
    m_{\rm mass}
    =
    \left[
        \frac18,\frac68,\frac18
    \right].
\end{equation}
%
Thus the mass contribution is
%
\begin{equation}
    \frac18u_{i-1}
    +\frac68u_i
    +\frac18u_{i+1}.
\end{equation}
%
Its Taylor expansion is
%
\begin{equation}
    \boxed{
    \frac18u_{i-1}
    +\frac68u_i
    +\frac18u_{i+1}
    =
    u_i
    +\frac18\Delta x^2u_{xx}
    +\frac1{96}\Delta x^4u_{xxxx}
    +O(\Delta x^6).
    }
\end{equation}
%
Combining this with the artificial-viscosity matrix gives
%
\begin{align}
    Au
    &=
    u
    +\left(\frac18-c_\psi\right)
      \Delta x^2u_{xx}
    \nonumber\\
    &\quad
    +\left(\frac1{96}-\frac{c_\psi}{12}\right)
      \Delta x^4u_{xxxx}
    +O(\Delta x^6).
\end{align}
%
Hence
%
\begin{equation}
    \boxed{
    Au=
    u+
    \left(\frac18-c_\psi\right)\Delta x^2u_{xx}
    +O(\Delta x^4).
    }
\end{equation}
%
This shows that the artificial-viscosity operator competes with the
second-derivative contribution already present in the mass stencil.
%
\section{Additional $c_E$ contribution}
%
The numerical scheme uses
%
\begin{equation}
    c_E=
    c_\psi^2\frac{\pi}{2}.
\end{equation}
%
The right-hand side is
%
\begin{equation}
    {\rm rhs}_i
    =
    c_E
    \left(
        \frac12\Delta x\,|u_{xx}|
    \right).
\end{equation}
%
Consequently,
%
\begin{equation}
    \boxed{
    {\rm rhs}_i=
    \frac{\pi}{4}
    c_\psi^2\Delta x\,|u_{xx}|.
    }
\end{equation}
%
This means that the complete artificial-viscosity mechanism is not
equivalent to a simple constant-viscosity Burgers equation. It is
solution dependent because it contains $|u_{xx}|$.

An effective viscosity estimate based on interpreting this correction
as a diffusive flux can be written schematically as
%
\begin{equation}
    \nu_{\rm eff}(x)
    \sim
    \frac{\pi}{4}
    c_\psi^2\Delta x
    \frac{|u_{xx}|}{|u_x|},
\end{equation}
%
where this relation is an interpretation rather than an exact
identification unless the complete time-stepping equation is specified.
%
\section{Shock-width estimate}
%
If one interprets the $c_E$ contribution as producing an effective
viscosity
%
\begin{equation}
    \nu_{\rm eff}
    \approx
    \frac{\pi}{4}c_\psi^2\Delta x,
\end{equation}
%
then for $a=1$ the characteristic shock length is approximately
%
\begin{equation}
    \delta
    \approx
    2\nu_{\rm eff}
    =
    \frac{\pi}{2}c_\psi^2\Delta x.
\end{equation}
%
In terms of grid cells,
%
\begin{equation}
    N_s=\frac{\delta}{\Delta x},
\end{equation}
%
so
%
\begin{equation}
    \boxed{
    N_s\approx\frac{\pi}{2}c_\psi^2.
    }
\end{equation}
%
Solving for $c_\psi$ gives
%
\begin{equation}
    \boxed{
    c_\psi\approx
    \sqrt{\frac{2N_s}{\pi}}.
    }
\end{equation}
%
Some example values are
%
\begin{center}
\begin{tabular}{cc}
\toprule
Desired shock width & $c_\psi$\\
\midrule
2 cells & 1.13\\
3 cells & 1.38\\
4 cells & 1.60\\
5 cells & 1.78\\
6 cells & 1.95\\
\bottomrule
\end{tabular}
\end{center}
%
Thus a reasonable starting value for a shock of about four grid cells
is
%
\begin{equation}
    \boxed{c_\psi\approx1.6}.
\end{equation}
%
These values are estimates because the exact effective viscosity
depends on how the matrix $A$ and the right-hand side are used in the
time-stepping equation.
%
\newpage
\section{Original C++ artificial-viscosity code}
%
The relevant code discussed was
%
{{\small
\begin{lstlisting}[language=C++]
double c_error = c_psi;
double c_E = c_psi * c_psi * std::numbers::pi / 2.0;

for (size_t i = 1; i < nx - 1; ++i)
{
    A.coeffRef(i, i - 1) = m_mass[0] - c_error;
    A.coeffRef(i, i    ) = m_mass[1] + 2. * c_error;
    A.coeffRef(i, i + 1) = m_mass[2] - c_error;

    rhs[i] = c_E * (0.5 * dx * std::abs(u_xixi[i]));
}
\end{lstlisting}
}}
%
with
%
\begin{equation}
    m_{\rm mass}
    =
    \left[
        \frac18,\frac68,\frac18
    \right].
\end{equation}
%
The exact effective viscosity cannot be determined uniquely from
these lines alone. The remaining required information is the complete
equation in which $A$ and ${\rm rhs}$ are used, particularly the
Eigen solve and the time-stepping terms.
%
\section{Recommended interpretation}
%
For a controlled numerical shock in Burgers' equation, a useful target
is a shock spanning approximately $3$--$5$ grid cells. The artificial
viscosity should be strong near the shock and weak in smooth regions.

A common form is
%
\begin{equation}
    u_t+u u_x
    =
    \frac{\partial}{\partial x}
    \left(
        \nu_{\rm art}u_x
    \right),
\end{equation}
%
with a solution-dependent viscosity such as
%
\begin{equation}
    \nu_{\rm art}
    =
    C_\nu\Delta x\,|u|
\end{equation}
%
or a shock detector based on the local gradient or curvature.
%
For the stationary shock
%
\[
a=1,\qquad b=-1,
\]
%
the ideal reference solution remains
%
\begin{equation}
    u(x)=
    -\tanh\left[
        \frac{x-x_s}{2\nu_{\rm art}}
    \right].
\end{equation}
%
The numerical solution can be compared against this profile to
quantify shock width and wiggles.
%
\section{Conclusion}
%
For
%
\[
u_t+uu_x=0,
\qquad
u(0,t)=a,
\qquad
u(L_x,t)=b,
\]
%
the inviscid shock speed is
%
\begin{equation}
    \boxed{s=\frac{a+b}{2}}.
\end{equation}
%
A stationary shock requires
%
\begin{equation}
    \boxed{b=-a}.
\end{equation}
%
With constant viscosity, the stationary viscous shock for $a>0$ and
$b=-a$ is
%
\begin{equation}
    \boxed{
    u(x)=
    -a\tanh\left[
        \frac{a(x-x_s)}{2\nu_{\rm art}}
    \right].
    }
\end{equation}
%
For the numerical scheme discussed here,
%
\[
m_{\rm mass}
=
\left[
\frac18,\frac68,\frac18
\right],
\]
%
and the artificial-viscosity stencil is
%
\[
c_\psi[-1,2,-1].
\]
%
Its leading Taylor contribution is proportional to
%
\[
-c_\psi\Delta x^2u_{xx},
\]
%
while the additional right-hand-side correction is
%
\[
\frac{\pi}{4}c_\psi^2\Delta x|u_{xx}|.
\]
%
The exact continuous equation represented by the complete numerical
scheme requires the remaining time-stepping/linear-system equations.

\end{document}
